% GENERATED FILE. Edit canonical lessons, then run npm run generate:books.
% canonical-source-hash: fnv1a64:d4d97939137c5648
% canonical-lessons: GU-C139-mokalvu, GU-C139-bhulvu, GU-C139-sodhvu, GU-C139-vaparvu, GU-C139-banavvu

\chapter{To send, To forget, To look for, To use, To make}
\label{ch:gu-to-send-to-forget-to-look-for-to-use-to-make}
\hlchaptermodality{139}

\begin{chapteropening}
\textbf{By the end of this chapter:} \emph{I can say to send, to forget, to look for, to use and to make.}

Complete the last lesson of chapter 139: I can say to send, to forget, to look for, to use and to make.
\end{chapteropening}

\section[\texorpdfstring{\gu{મોકલવું} (mokalvũ)}{મોકલવું (mokalvũ)}]{\gu{મોકલવું} (mokalvũ) — to send}

\label{lesson:GU-C139-mokalvu}



\begin{quote}
Before the new one: say the Gujarati for to open, then the Gujarati for to call.
\end{quote}

\subsection*{You'll want to know: \gu{મોકલવું}}
\textbf{\gu{મોકલવું}} (\emph{mokalvũ}) — \textquotedblleft{}to send\textquotedblright{}.

\begin{grammarlens}[title={What you've built}]
One more word for this chapter.
\end{grammarlens}

\subsection*{Your turn}
\begin{itemize}
\raggedright
  \item \emph{Say it:} \emph{mokalvũ}
  \item \emph{Say it:} \emph{mokalvũ}, once more
  \item \emph{Say it:} say \emph{mokalvũ}
  \item \emph{Recall:} say the Gujarati for to open, then the Gujarati for to call, then say \emph{mokalvũ} again
\end{itemize}

\begin{tcolorbox}[breakable,colback=teal!4,colframe=teal!35!black,title={Before you move on}]
What does \gu{મોકલવું} mean? (\textquotedblleft{}To send\textquotedblright{}.) Say it once more.
\end{tcolorbox}
\section[\texorpdfstring{\gu{ભૂલવું} (bhūlvũ)}{ભૂલવું (bhūlvũ)}]{\gu{ભૂલવું} (bhūlvũ) — to forget}

\label{lesson:GU-C139-bhulvu}



\begin{quote}
Before the new one: say the Gujarati for to call, then the Gujarati for to send.
\end{quote}

\subsection*{You'll want to know: \gu{ભૂલવું}}
\textbf{\gu{ભૂલવું}} (\emph{bhūlvũ}) — \textquotedblleft{}to forget\textquotedblright{}.

\begin{grammarlens}[title={What you've built}]
One more word for this chapter.
\end{grammarlens}

\subsection*{Your turn}
\begin{itemize}
\raggedright
  \item \emph{Say it:} \emph{bhūlvũ}
  \item \emph{Say it:} \emph{bhūlvũ}, once more
  \item \emph{Say it:} say \emph{bhūlvũ}
  \item \emph{Recall:} say the Gujarati for to call, then the Gujarati for to send, then say \emph{bhūlvũ} again
\end{itemize}

\begin{tcolorbox}[breakable,colback=teal!4,colframe=teal!35!black,title={Before you move on}]
What does \gu{ભૂલવું} mean? (\textquotedblleft{}To forget\textquotedblright{}.) Say it once more.
\end{tcolorbox}
\section[\texorpdfstring{\gu{શોધવું} (śodhvũ)}{શોધવું (śodhvũ)}]{\gu{શોધવું} (śodhvũ) — to look for}

\label{lesson:GU-C139-sodhvu}



\begin{quote}
Before the new one: say the Gujarati for to send, then the Gujarati for to forget.
\end{quote}

\subsection*{You'll want to know: \gu{શોધવું}}
\textbf{\gu{શોધવું}} (\emph{śodhvũ}) — \textquotedblleft{}to look for\textquotedblright{}.

\begin{grammarlens}[title={What you've built}]
One more word for this chapter.
\end{grammarlens}

\subsection*{Your turn}
\begin{itemize}
\raggedright
  \item \emph{Say it:} \emph{śodhvũ}
  \item \emph{Say it:} \emph{śodhvũ}, once more
  \item \emph{Say it:} say \emph{śodhvũ}
  \item \emph{Recall:} say the Gujarati for to send, then the Gujarati for to forget, then say \emph{śodhvũ} again
\end{itemize}

\begin{tcolorbox}[breakable,colback=teal!4,colframe=teal!35!black,title={Before you move on}]
What does \gu{શોધવું} mean? (\textquotedblleft{}To look for\textquotedblright{}.) Say it once more.
\end{tcolorbox}
\section[\texorpdfstring{\gu{વાપરવું} (vāparvũ)}{વાપરવું (vāparvũ)}]{\gu{વાપરવું} (vāparvũ) — to use}

\label{lesson:GU-C139-vaparvu}



\begin{quote}
Before the new one: say the Gujarati for to forget, then the Gujarati for to look for.
\end{quote}

\subsection*{You'll want to know: \gu{વાપરવું}}
\textbf{\gu{વાપરવું}} (\emph{vāparvũ}) — \textquotedblleft{}to use\textquotedblright{}.

\begin{grammarlens}[title={What you've built}]
One more word for this chapter.
\end{grammarlens}

\subsection*{Your turn}
\begin{itemize}
\raggedright
  \item \emph{Say it:} \emph{vāparvũ}
  \item \emph{Say it:} \emph{vāparvũ}, once more
  \item \emph{Say it:} say \emph{vāparvũ}
  \item \emph{Recall:} say the Gujarati for to forget, then the Gujarati for to look for, then say \emph{vāparvũ} again
\end{itemize}

\begin{tcolorbox}[breakable,colback=teal!4,colframe=teal!35!black,title={Before you move on}]
What does \gu{વાપરવું} mean? (\textquotedblleft{}To use\textquotedblright{}.) Say it once more.
\end{tcolorbox}
\section[\texorpdfstring{\gu{બનાવવું} (banāvvũ)}{બનાવવું (banāvvũ)}]{\gu{બનાવવું} (banāvvũ) — to make}

\label{lesson:GU-C139-banavvu}



\begin{quote}
Before the new one: say the Gujarati for to look for, then the Gujarati for to use.
\end{quote}

\subsection*{You'll want to know: \gu{બનાવવું}}
\textbf{\gu{બનાવવું}} (\emph{banāvvũ}) — \textquotedblleft{}to make\textquotedblright{}.

\begin{grammarlens}[title={What you've built}]
One more word for this chapter.
\end{grammarlens}

\subsection*{Your turn}
\begin{itemize}
\raggedright
  \item \emph{Say it:} \emph{banāvvũ}
  \item \emph{Say it:} \emph{banāvvũ}, once more
  \item \emph{Say it:} say \emph{banāvvũ}
  \item \emph{Recall:} say the Gujarati for to look for, then the Gujarati for to use, then say \emph{banāvvũ} again
\end{itemize}

\begin{tcolorbox}[breakable,colback=teal!4,colframe=teal!35!black,title={Before you move on}]
What does \gu{બનાવવું} mean? (\textquotedblleft{}To make\textquotedblright{}.) Say it once more.
\end{tcolorbox}
